% ============================================================================
% Chapter 7: κ = 2π from Homotopy
% ============================================================================
% Source: DERIVE_KAPPA_FROM_5D_HOMOTOPY.md
% Status: FILLED from SoT
% Contamination check: PASSED (no SM terms)
% ============================================================================

\chapter{$\kappa = 2\pi$ from Homotopy}
\label{ch:kappa}

\source{DERIVE\_KAPPA\_FROM\_5D\_HOMOTOPY.md}

\begin{abstract}
Chapter~\ref{ch:instanton} established that the Euclidean action for metastable junction
decay has the form $\SE/\hbar = \kappa \times (L_0/\delta)$. What fixes $\kappa$?
This chapter derives $\kappa = 2\pi$ from the fundamental group of the circle,
$\pione(\Sone) = \mathbb{Z}$. The result is topological: $\kappa$ is not a fit
parameter but a consequence of winding number quantization in the compact
fifth dimension.
\end{abstract}

% ============================================================================
\section{The Winding Number}
\label{sec:kappa:winding}

\subsection{Topological Invariant}

In the EDC framework, the metastable junction is a junction defect with structure in the
fifth dimension. Decay corresponds to a \emph{topological transition}---a change
in the configuration's winding around the compact direction.

\begin{definition}[Winding Number]
    Let $\theta \in [0, 2\pi)$ parametrize the compact fifth dimension (with
    topology $\Sone$). A field configuration $\phi$ has \textbf{winding number}
    $w \in \mathbb{Z}$ if:
    \begin{equation}
        \phi(\theta + 2\pi) = \phi(\theta) + 2\pi w
        \label{eq:winding_bc}
    \end{equation}
    The integer $w$ counts how many times the configuration ``wraps'' around
    the compact direction. \tagDer
\end{definition}

\textbf{Key property:} The winding number is a \emph{topological invariant}.
It cannot change under continuous deformations---only discrete transitions
(such as tunneling) can alter $w$.

\subsection{Winding and the Instanton}

An instanton is a Euclidean-time configuration that interpolates between
topologically distinct states. For the compact dimension with $\Sone$ topology:
\begin{itemize}
    \item Initial state: winding number $w$
    \item Final state: winding number $w \pm 1$
    \item Instanton: mediates the transition $\Delta w = \pm 1$
\end{itemize}

The \emph{minimal} instanton has $|\Delta w| = 1$. Higher winding transitions
($|\Delta w| > 1$) are exponentially suppressed relative to the fundamental one.

% ============================================================================
\section{$\pione(\Sone) = \mathbb{Z}$}
\label{sec:kappa:fundamental}

\subsection{The Fundamental Group}

The fundamental group $\pione(X)$ of a topological space $X$ classifies loops
in $X$ up to continuous deformation (homotopy).

\begin{derivationbox}[title=Fundamental Group of the Circle]
    For the circle $\Sone$:
    \begin{equation}
        \boxed{\pione(\Sone) = \mathbb{Z}}
        \label{eq:pi1_S1}
    \end{equation}

    \textbf{Proof sketch:}
    \begin{enumerate}
        \item A loop $\gamma: [0,1] \to \Sone$ with $\gamma(0) = \gamma(1) = p_0$
              can be lifted to the universal cover $\mathbb{R}$.
        \item The lift $\tilde{\gamma}: [0,1] \to \mathbb{R}$ satisfies
              $\tilde{\gamma}(0) = 0$ and $\tilde{\gamma}(1) = 2\pi n$
              for some $n \in \mathbb{Z}$.
        \item The integer $n$ is the winding number of the loop.
        \item Two loops are homotopic if and only if they have the same $n$.
    \end{enumerate}
    \tagDer
\end{derivationbox}

\subsection{Interpretation}

Each element $n \in \mathbb{Z}$ corresponds to a homotopy class of loops:
\begin{center}
\begin{tabular}{cl}
\toprule
$n$ & Homotopy class \\
\midrule
$0$ & Contractible loop (trivial) \\
$+1$ & Single counterclockwise wrap \\
$-1$ & Single clockwise wrap \\
$+k$ & $k$ counterclockwise wraps \\
\bottomrule
\end{tabular}
\end{center}

The group operation is loop concatenation: $n_1 + n_2$ wraps.

% ============================================================================
\section{Physical Winding in EDC}
\label{sec:kappa:physical}

\subsection{The Compact Dimension}

In EDC, the brane extends in four large dimensions plus one compact direction.
The compact dimension has characteristic scale $\delta$ and topology $\Sone$.

\begin{edcopenbox}[title=Assumption: $\Sone$ Topology]
    \textbf{Assumption:} The compact fifth dimension relevant to junction
    dynamics has the topology of a circle $\Sone$.

    This is a structural assumption [P], not derived from more fundamental
    principles within the current framework. \tagP
\end{edcopenbox}

\subsection{Junction States and Winding}

The anchor and metastable junction are interpreted as junction configurations differing
in their topological class:

\begin{center}
\begin{tabular}{lll}
\toprule
Configuration & Winding & Energy \\
\midrule
Anchor & $w = 0$ & Ground state \\
Metastable & $w = 1$ & Metastable (higher by $\Delta m_{np}$) \\
\bottomrule
\end{tabular}
\end{center}

\textbf{Decay as unwinding:} Metastable decay corresponds to a topological
transition from $w = 1$ to $w = 0$:
\begin{equation}
    \Delta w = -1 \quad \text{(unwinding by one unit)}
    \label{eq:delta_w}
    \tagP
\end{equation}

This transition cannot occur classically (topology is preserved under
continuous evolution) but can proceed via quantum tunneling---the instanton.

\subsection{Minimal Transition}

The fundamental instanton has $|\Delta w| = 1$. This is the dominant decay
channel because:
\begin{enumerate}
    \item Higher winding transitions require traversing multiple topological
          sectors
    \item The action scales with $|\Delta w|$, making $|\Delta w| > 1$
          transitions exponentially suppressed
\end{enumerate}

For metastable decay: the instanton has $\Delta w = -1$ (or equivalently, the
Euclidean configuration has winding $w = 1$ in imaginary time).

% ============================================================================
\section{Derivation: $\kappa = 2\pi$}
\label{sec:kappa:result}

\subsection{The Winding Integral}

The topological contribution to the Euclidean action arises from the angular
integration around the compact dimension.

\begin{derivationbox}[title=Angular Integration]
    Let $\theta$ parametrize $\Sone$ with period $2\pi$. For a configuration
    with winding number $w$, the angular gradient is:
    \begin{equation}
        \frac{\partial \phi}{\partial \theta} = \frac{w}{R}
    \end{equation}
    where $R$ is the radius of the compact dimension.

    The integral around the compact direction gives:
    \begin{equation}
        \oint_{\Sone} d\theta = 2\pi
        \label{eq:angular_integral}
    \end{equation}

    For a winding-$w$ configuration, the topological charge accumulated is:
    \begin{equation}
        Q_{\text{top}} = \frac{1}{2\pi} \oint d\theta \, \frac{\partial\phi}{\partial\theta}
                       = \frac{1}{2\pi} \times 2\pi \times w = w
    \end{equation}
    \tagDer
\end{derivationbox}

\subsection{Action Normalization}

The Euclidean action for a topological transition with charge $Q_{\text{top}} = w$
has the canonical form:
\begin{equation}
    \SE = 2\pi \, |w| \times (\text{geometric scale factor})
    \label{eq:SE_canonical}
    \tagDc
\end{equation}

\textbf{Why the factor $2\pi$?}

The $2\pi$ arises from three equivalent perspectives:
\begin{enumerate}
    \item \textbf{Angular measure:} $\oint d\theta = 2\pi$ is the circumference
          of the unit circle in radians.
    \item \textbf{Flux quantization:} The ``flux'' through a winding-$w$
          configuration is $2\pi w$, giving action $\propto 2\pi |w|$.
    \item \textbf{Topological normalization:} The action must be $2\pi \times$
          (integer) for the path integral to be single-valued under large
          coordinate transformations on $\Sone$.
\end{enumerate}

\subsection{The EDC Result}

In the EDC framework, the geometric scale factor is the dimensionless ratio
$\Lzero/\bdelta$:

\begin{resultbox}[title=Topological Factor $\kappa$]
    For a topological transition with winding change $|\Delta w| = |w|$:
    \begin{equation}
        \frac{\SE}{\hbar} = 2\pi \, |w| \times \frac{\Lzero}{\bdelta}
        \label{eq:SE_full}
    \end{equation}

    For the fundamental instanton ($|w| = 1$):
    \begin{equation}
        \boxed{\kappa = 2\pi}
        \label{eq:kappa_result}
    \end{equation}
    \tagDc
\end{resultbox}

\textbf{Epistemic status:} The result $\kappa = 2\pi$ is [Dc]---derived with
the constraint that the compact dimension has $\Sone$ topology [P]. Given this
assumption, $\kappa = 2\pi$ follows from standard topology ($\pione(\Sone) = \mathbb{Z}$)
and canonical action normalization.

\subsection{What $\kappa = 2\pi$ Is NOT}

To clarify the nature of this result:
\begin{itemize}
    \item $\kappa$ is \emph{not} a fit parameter adjusted to match data
    \item $\kappa$ is \emph{not} determined by dynamics or energy scales
    \item $\kappa$ \emph{is} fixed by topology: once $\Sone$ is assumed,
          $\kappa = 2\pi$ is forced
\end{itemize}

The only freedom is whether the compact dimension has $\Sone$ topology (vs.\
some other manifold). Different topologies would give different values of
$\kappa$---see the falsifiability discussion below.

% ============================================================================
\section{Epistemic Status}
\label{sec:kappa:epistemic}

\begin{center}
\begin{tabular}{llll}
\toprule
\textbf{Claim} & \textbf{Status} & \textbf{Reference} & \textbf{Dependency} \\
\midrule
$\pione(\Sone) = \mathbb{Z}$ & [Der] & Eq.~\eqref{eq:pi1_S1} & Standard topology \\
$\oint d\theta = 2\pi$ & [Der] & Eq.~\eqref{eq:angular_integral} & Definition \\
Action $\propto 2\pi |w|$ & [Dc] & Eq.~\eqref{eq:SE_canonical} & Normalization \\
$\kappa = 2\pi$ & [Dc] & Eq.~\eqref{eq:kappa_result} & $\Sone$ topology [P] \\
\midrule
Compact dim.\ has $\Sone$ & [P] & \S\ref{sec:kappa:physical} & Assumption \\
Metastable has $w = 1$ & [P] & \S\ref{sec:kappa:physical} & Model interpretation \\
Decay is $\Delta w = -1$ & [P] & Eq.~\eqref{eq:delta_w} & Model interpretation \\
\bottomrule
\end{tabular}
\end{center}

% ============================================================================
\section{Falsifiability}
\label{sec:kappa:falsify}

\begin{edcopenbox}[title=Assumption Check and Falsifiability]
    \textbf{Central assumption:} The compact dimension relevant to instanton
    tunneling has $\Sone$ topology. \tagP

    \textbf{What would falsify this:}
    \begin{enumerate}
        \item \textbf{Different topology:} If the compact dimension were
              $\Sone \times \Sone$ (torus), the relevant homotopy group would
              be $\pione(T^2) = \mathbb{Z} \times \mathbb{Z}$, potentially
              giving a different $\kappa$.

        \item \textbf{Trivial $\pione$:} If $\pione = 0$ (simply connected
              compact space), there would be no topological protection---no
              winding, no instanton, and the decay mechanism would require
              revision.

        \item \textbf{Experimental mismatch:} If the full formula
              $\taun = A \cdot (\hbar/\omegaz) \cdot \exp[\kappa(L_0/\delta)]$
              with $\kappa = 2\pi$ cannot match $\taun \approx 880\second$
              for reasonable values of other parameters, the $\Sone$ assumption
              would be suspect.
    \end{enumerate}

    \textbf{Current status:} The assumption is consistent with the observed
    metastable lifetime when combined with $L_0/\delta \approx \pi^2$
    (Chapter~\ref{ch:L0delta}).
\end{edcopenbox}

% ============================================================================
% Observer-side projection
% ============================================================================

\begin{observerbox}
Observer-side projection note: quoted terms are measurement labels for the
3D/4D projection of the 5D objects defined in this chapter; they do not
imply any conventional mechanism.

\begin{itemize}
    \item \MetastableJunction{} $\leftrightarrow$ ``neutron'' (projection label)
\end{itemize}

What the observer records: the topological factor $\kappa = 2\pi$ enters
the Euclidean action formula and directly affects the measured lifetime
$\taun$ of the metastable species at the projection level.
\end{observerbox}

% ============================================================================
\section{Summary}
\label{sec:kappa:summary}

\begin{resultbox}
    The topological factor $\kappa = 2\pi$ is derived from first principles:
    \begin{enumerate}
        \item The compact fifth dimension has $\Sone$ topology [P]
        \item The fundamental group is $\pione(\Sone) = \mathbb{Z}$ [Der]
        \item Instanton transitions carry winding $|w| = 1$ [P]
        \item Angular integration gives $\oint d\theta = 2\pi$ [Der]
        \item Therefore: $\kappa = 2\pi$ [Dc]
    \end{enumerate}

    \textbf{Key point:} $\kappa$ is not a fit parameter. It is fixed by topology
    once the $\Sone$ structure is assumed.
\end{resultbox}

\bigskip
\noindent
\textbf{Forward references:}
\begin{itemize}
    \item Chapter~\ref{ch:L0delta} derives the geometric ratio $\Lzero/\bdelta$
    \item Chapter~\ref{ch:tau_n} combines $\kappa$ and $\Lzero/\bdelta$ to
          compute $\taun \approx 880\second$
\end{itemize}

\bigskip
\noindent
\textbf{Derivation chain status:}
\begin{equation*}
    \SE/\hbar = \underbrace{2\pi}_{\kappa \text{ [Dc, this chapter]}}
              \times \underbrace{L_0/\delta}_{\text{[P], Ch.~\ref{ch:L0delta}}}
\end{equation*}

